\subsection{Matemática}
\subsubsection{Constantes}

\begin{lstlisting}[language=c++, breaklines=true]
/* Definitions of useful mathematical constants
 * M_E        - e
 * M_LOG2E    - log2(e)
 * M_LOG10E   - log10(e)
 * M_LN2      - ln(2)
 * M_LN10     - ln(10)
 * M_PI       - pi
 * M_PI_2     - pi/2
 * M_PI_4     - pi/4
 * M_1_PI     - 1/pi
 * M_2_PI     - 2/pi
 * M_2_SQRTPI - 2/sqrt(pi)
 * M_SQRT2    - sqrt(2)
 * M_SQRT1_2  - 1/sqrt(2)
 */

#define M_E        2.71828182845904523536
#define M_LOG2E    1.44269504088896340736
#define M_LOG10E   0.434294481903251827651
#define M_LN2      0.693147180559945309417
#define M_LN10     2.30258509299404568402
#define M_PI       3.14159265358979323846
#define M_PI_2     1.57079632679489661923
#define M_PI_4     0.785398163397448309616
#define M_1_PI     0.318309886183790671538
#define M_2_PI     0.636619772367581343076
#define M_2_SQRTPI 1.12837916709551257390
#define M_SQRT2    1.41421356237309504880
#define M_SQRT1_2  0.707106781186547524401
\end{lstlisting}

\subsubsection{Fórmulas e Teoremas}

\textbf{Teorema de Lagrange:} Todo inteiro positivo pode ser representado como soma de $4$ quadrados.

\textbf{Teorema 'Eureka' de Gauss:} Todo inteiro positivo pode ser expresso como a soma de três números triangulares.

\textbf{Teorema de Zeckendorf:} Todo inteiro positivo tem uma representação única como soma de números de Fibonacci tal que não existem dois números iguais ou consecutivos.

\textbf{Fórmula de Euclides para ternas pitagóricas:} Ternas pitagóricas primitivas são todas da forma 
$$ (n^2 - m^2, 2nm, n^2 + m^2)$$
em que $0 < m < n$, $(n, m) = 1$ e $2 | nm$.
\subsubsection{Números primos}
